% ============================================================================
\section{Paradigm 4: Kolmogorov-Arnold Networks and LUT Compilation}
\label{sec:kan}
% ============================================================================

% --- Intuition ---
\subsection{The Idea in Plain English}

The three paradigms we have seen so far all simplify the neuron: strip
it down to binary operations and accept some accuracy loss. This
fourth paradigm does the opposite during training: it makes each
connection \emph{more} complex than a traditional weight. Then, once
training is done, it compiles that complexity into a simple lookup
table.

\begin{analogy}
Imagine learning to cook a complex sauce. At first, you follow an
elaborate recipe (training): you measure precisely, taste repeatedly,
adjust seasoning. Once you have perfected it, you write down just
the final quantities on an index card (lookup table). Now you can
recreate the sauce by just reading the card---no skill required,
just follow the numbers. The complexity happened once; the
deployment is trivial.
\end{analogy}

The ``elaborate recipe'' is a \concept{Kolmogorov-Arnold Network}
(KAN)~\cite{liu2024kan}, where every connection is a learnable curve (B-spline). The
``index card'' is a \concept{lookup table} (LUT), where the curve is
sampled at fixed points and stored in memory.

% --- Mechanism ---
\subsection{The Kolmogorov-Arnold Theorem}

KANs are grounded in a mathematical theorem from 1957~\cite{kolmogorov1957}. In plain
terms:

\begin{keyinsight}[title={Kolmogorov-Arnold Representation Theorem~\cite{kolmogorov1957}}]
\emph{Any} continuous function of multiple variables can be written as
a composition of continuous functions of \emph{one} variable and
addition. You never need multiplication of two variables---just
clever single-variable functions and sums.
\end{keyinsight}

Formally, for any continuous function $f : [0,1]^n \to \mathbb{R}$:
\begin{equation}
  f(x_1, \ldots, x_n) = \sum_{q=0}^{2n}
    \Phi_q\!\left(\sum_{p=1}^{n} \phi_{q,p}(x_p)\right)
  \label{eq:ka-theorem}
\end{equation}

where $\phi_{q,p}$ and $\Phi_q$ are continuous single-variable
functions. The theorem says these functions exist but does not tell
us what they look like---the network learns them.

\subsection{How a KAN Layer Works}

\subsubsection{Traditional MLP: Fixed Activations, Learnable Weights}

Recall that in a standard MLP, each connection has a single number
(weight), and the activation function is the same for all neurons in
a layer:

\begin{equation}
  \text{MLP layer: } \bm{y} = \sigma(\bm{W}\bm{x} + \bm{b})
  \label{eq:mlp-layer}
\end{equation}

\subsubsection{KAN: Learnable Activations, No Weights}

A KAN~\cite{liu2024kan} \emph{replaces} the weight-then-activate pattern. Each edge has
its own learnable function $\phi_{i,j}$. Nodes just sum:

\begin{equation}
  \text{KAN layer: } y_i = \sum_{j=1}^{n} \phi_{i,j}(x_j)
  \label{eq:kan-layer}
\end{equation}

\begin{figure}[htbp]
\centering
\begin{tikzpicture}[
  node/.style={draw, circle, minimum size=0.7cm, font=\small, inner sep=0pt},
  >=Stealth
]
  % --- MLP side ---
  \node[font=\sffamily\bfseries] at (2, 4.5) {MLP};

  % Input
  \foreach \i in {1,2,3} {
    \node[node, fill=lightbg] (mi\i) at (0, 3-\i) {$x_\i$};
  }
  % Output
  \foreach \i in {1,2} {
    \node[node, fill=bitblue!15] (mo\i) at (4, 2.5-\i) {$\sigma$};
  }
  % Connections (straight lines = scalar weights)
  \foreach \i in {1,2,3} {
    \foreach \j in {1,2} {
      \draw[->, thick, bitgray] (mi\i) -- (mo\j)
        node[midway, font=\tiny, fill=white, inner sep=1pt] {$w$};
    }
  }

  % --- KAN side ---
  \node[font=\sffamily\bfseries] at (9, 4.5) {KAN};

  % Input
  \foreach \i in {1,2,3} {
    \node[node, fill=lightbg] (ki\i) at (7, 3-\i) {$x_\i$};
  }
  % Output
  \foreach \i in {1,2} {
    \node[node, fill=bitorange!15] (ko\i) at (11, 2.5-\i) {$\Sigma$};
  }
  % Connections (wavy lines = learnable functions)
  \foreach \i in {1,2,3} {
    \foreach \j in {1,2} {
      \draw[->, thick, bitorange, decorate,
        decoration={snake, amplitude=1.5pt, segment length=8pt}]
        (ki\i) -- (ko\j)
        node[midway, font=\tiny, fill=white, inner sep=1pt] {$\phi$};
    }
  }

  % Labels
  \node[font=\scriptsize\sffamily, bitgray, align=center] at (2, -1.5)
    {Scalar weights on edges\\Fixed activation on nodes};
  \node[font=\scriptsize\sffamily, bitgray, align=center] at (9, -1.5)
    {Learnable functions on edges\\Simple summation on nodes};
\end{tikzpicture}
\caption{MLP vs KAN architecture. In an MLP (left), edges carry scalar
  weights and nodes apply a fixed activation. In a KAN (right), edges
  carry learnable functions (B-splines) and nodes simply sum. The wavy
  lines represent that each connection is an entire curve, not a single
  number.}
\label{fig:mlp-vs-kan}
\end{figure}

\subsubsection{B-Spline Functions}

Each edge function $\phi_{i,j}$ is parameterised as a B-spline with
$G$ grid points and order $k$ (typically $k=3$, cubic). The function
is:
\begin{equation}
  \phi_{i,j}(x) = w_b \cdot \operatorname{SiLU}(x)
    + w_s \cdot \sum_{m=0}^{G+k} c_m \cdot B_m^k(x)
  \label{eq:kan-spline}
\end{equation}

where $B_m^k$ are B-spline basis functions, $c_m$ are learnable
coefficients, and $w_b, w_s$ are scaling factors.

\begin{engineernote}
A B-spline is just a smooth, piecewise polynomial curve controlled
by a set of ``control points'' (the coefficients $c_m$). If you have
used Bézier curves in graphics, B-splines are their generalisation.
The grid determines where the curve's ``knots'' are---points where
one polynomial piece ends and the next begins.
\end{engineernote}

\begin{figure}[htbp]
\centering
\begin{tikzpicture}
  \begin{axis}[
    width=10cm, height=5cm,
    xlabel={Input $x$}, ylabel={$\phi(x)$},
    grid=major, grid style={bitgray!20},
    xmin=-2, xmax=2, ymin=-1.5, ymax=2,
    axis lines=middle,
    every axis label/.style={font=\small\sffamily},
    tick label style={font=\tiny}
  ]
    % A sample spline-like function
    \addplot[thick, bitblue, smooth, domain=-2:2, samples=100]
      {0.3*x + 0.5*sin(deg(2*x)) + 0.2*cos(deg(3*x))};
    \addlegendentry{Learned $\phi_{i,j}(x)$}

    % Grid points
    \foreach \gp in {-2, -1.5, -1, -0.5, 0, 0.5, 1, 1.5, 2} {
      \addplot[only marks, mark=|, mark size=3pt, bitgray]
        coordinates {(\gp, -1.5)};
    }

    % ReLU for comparison
    \addplot[thick, bitorange, dashed, domain=-2:2, samples=50]
      {max(0, x)};
    \addlegendentry{Fixed ReLU (MLP)}
  \end{axis}

  % Grid label
  \node[font=\scriptsize\sffamily, bitgray] at (1, -0.3)
    {Grid points (knots)};
\end{tikzpicture}
\caption{A learnable B-spline activation (blue) vs.\ a fixed ReLU
  (orange dashed). The B-spline can learn arbitrary shapes, capturing
  complex input--output relationships that a ReLU cannot. The grid
  points (tick marks at bottom) define the spline's resolution.}
\label{fig:kan-spline}
\end{figure}

% --- LUT Compilation ---
\subsection{From Splines to Lookup Tables}
\label{sec:lut-compilation}

Here is where the computation-memory trade-off becomes concrete. After
training, each spline function $\phi_{i,j}(x)$ is a fixed curve. We
can \emph{compile} it:

\begin{enumerate}[leftmargin=2em]
  \item \textbf{Sample} the spline at $L$ evenly spaced points.
  \item \textbf{Quantise} the output values to a fixed-point format
    (e.g., int8).
  \item \textbf{Store} the $L$ values in an array---the lookup table.
  \item At inference, \textbf{look up} the nearest entry (or
    interpolate between two adjacent entries).
\end{enumerate}

\begin{figure}[htbp]
\centering
\begin{tikzpicture}[>=Stealth]
  % Spline
  \begin{scope}
    \node[font=\small\sffamily\bfseries] at (2.5, 3.5) {Trained Spline};
    \draw[thick, bitblue, smooth] plot coordinates
      {(0,1) (0.5,1.8) (1,2.5) (1.5,2.2) (2,1.5) (2.5,0.8)
       (3,1.2) (3.5,2.0) (4,2.3) (4.5,1.9) (5,1.0)};
    \draw[bitgray!50] (0, 0) -- (5, 0);
    \draw[bitgray!50] (0, 0) -- (0, 3);
  \end{scope}

  % Arrow
  \draw[->, very thick, bitgray] (5.5, 1.5) -- (7, 1.5)
    node[midway, above, font=\scriptsize\sffamily] {sample \&};
  \node[font=\scriptsize\sffamily, bitgray] at (6.25, 1.1) {quantise};

  % LUT
  \begin{scope}[xshift=7.5cm]
    \node[font=\small\sffamily\bfseries] at (2, 3.5) {Lookup Table};
    \foreach \i/\v/\q in {0/1.0/25, 1/1.8/45, 2/2.5/63,
      3/2.2/56, 4/1.5/38, 5/0.8/20, 6/1.2/30, 7/2.0/51,
      8/2.3/58, 9/1.9/48, 10/1.0/25} {
      \pgfmathsetmacro{\ypos}{2.8 - \i * 0.27}
      \node[draw, minimum width=0.8cm, minimum height=0.25cm,
        font=\tiny\ttfamily, fill=bitblue!10, inner sep=0pt]
        at (1.5, \ypos) {\q};
      \node[font=\tiny\ttfamily, bitgray, left] at (0.9, \ypos) {[\i]};
    }
    \node[font=\scriptsize\sffamily, bitgray] at (2, -0.4)
      {int8 values};
  \end{scope}

  % Usage arrow
  \draw[->, thick, bitgreen] (11, 2) -- (12.5, 2);
  \node[font=\scriptsize\sffamily, bitgreen] at (11.75, 2.4) {index};

  % Inference
  \begin{scope}[xshift=13cm]
    \node[font=\small\sffamily\bfseries] at (1, 3.5) {Inference};
    \node[draw, rounded corners, fill=bitgreen!10,
      font=\scriptsize\sffamily, align=center, minimum width=2cm]
      at (1, 2) {Input $x$\\$\downarrow$\\table[\texttt{idx}]\\$\downarrow$\\Output};
    \node[font=\scriptsize\sffamily, bitgray, align=center] at (1, -0.2)
      {One memory\\read};
  \end{scope}
\end{tikzpicture}
\caption{LUT compilation: a trained B-spline (left) is sampled at
  regular intervals and quantised to int8 values (centre). At
  inference, evaluating the function becomes a single array index
  operation (right)---no arithmetic needed.}
\label{fig:lut-compilation}
\end{figure}

\begin{keyinsight}
After LUT compilation, the entire ``activation function'' is a single
memory read. The spline's mathematical complexity exists only during
training. At deployment, every edge in the network costs exactly one
table lookup---a cost comparable to a single bitwise operation on
most hardware.
\end{keyinsight}

% --- ULEEN ---
\subsection{ULEEN: The Extreme Case}
\label{sec:uleen}

The \concept{Ultra-Low-Energy Edge Neural Network}
(ULEEN)~\cite{uleen2023} takes the LUT idea to its logical extreme:
the entire network is just tables. No splines, no arithmetic, not even
during training. The input features are grouped, each group indexes
into a table, and the outputs are summed.

\begin{table}[htbp]
\centering
\caption{ULEEN results on MNIST.}
\label{tab:uleen}
\begin{tabular}{@{}lr@{}}
  \toprule
  \textbf{Metric} & \textbf{Value} \\
  \midrule
  Accuracy          & 96.2\% \\
  Model size         & 16.9\,kB \\
  Inference latency  & 0.21\,$\mu$s \\
  Throughput         & 14.3\,M inferences/s \\
  \bottomrule
\end{tabular}
\end{table}

% --- Formalisation ---
\subsection{Mathematical Summary}
\label{sec:kan-math}

A KAN layer $\Phi$ mapping $\mathbb{R}^{n_\text{in}}$ to
$\mathbb{R}^{n_\text{out}}$ is defined by a matrix of univariate
functions:
\begin{equation}
  \Phi = \{\phi_{i,j}\}_{1 \leq i \leq n_\text{out},\;
    1 \leq j \leq n_\text{in}}
  \label{eq:kan-matrix}
\end{equation}

The layer computation is:
\begin{equation}
  y_i = \sum_{j=1}^{n_\text{in}} \phi_{i,j}(x_j),
  \quad i = 1, \ldots, n_\text{out}
  \label{eq:kan-layer-formal}
\end{equation}

Each function has $G + k + 2$ learnable parameters (B-spline
coefficients plus scaling), compared to 1 parameter per edge in an MLP.
A KAN layer with width $n$ and grid size $G$ has $O(n^2 (G+k))$
parameters vs.\ $O(n^2)$ for an MLP layer.

After LUT compilation with $L$ samples per function, each
$\phi_{i,j}$ is replaced by:
\begin{equation}
  \hat{\phi}_{i,j}(x) = \operatorname{table}_{i,j}
    \!\left[\operatorname{round}\!\left(\frac{x - x_{\min}}
    {x_{\max} - x_{\min}} \cdot (L-1)\right)\right]
  \label{eq:lut-eval}
\end{equation}

This is an $O(1)$ memory access, reducing inference to
$n_\text{in} \times n_\text{out}$ table lookups plus
$n_\text{out}$ integer additions per layer.

\subsection{Strengths and Limitations}

\begin{description}[leftmargin=1em]
  \item[Strengths:]
    Theoretically grounded (Kolmogorov-Arnold theorem~\cite{kolmogorov1957}), highest
    accuracy among bitwise-deployable approaches, LUT compilation
    enables zero-arithmetic inference, naturally maps to FPGA LUT
    fabric~\cite{polylut2023, lutkan2025}, graceful accuracy-memory trade-off (more table entries
    = more accuracy).

  \item[Limitations:]
    Training is slow ($\sim$10$\times$ slower than MLP with same
    parameter count), requires floating-point during training,
    LUT size grows with input precision and grid resolution,
    relatively new (2024--2025) with evolving best practices.
\end{description}
